% SOURCE_STATEMENT: "Can one establish local existence and uniqueness of solutions, together with continuous dependence on initial data, for the Lovelock-theory equations?"
% Preparation and final self-audit: CONTRIBUTING.md prompts C and D; source/audit artifacts retained in the local physics_statements_2026_10_11 batch.
% !TeX program = lualatex
\documentclass[11pt]{article}
\usepackage[margin=25mm]{geometry}
\usepackage{fontspec}
\setmainfont{Latin Modern Roman}
\usepackage{amsmath,amssymb,mathtools}
\usepackage{unicode-math}
\setmathfont{Latin Modern Math}
\usepackage{xurl,hyperref}
\hypersetup{colorlinks=true,urlcolor=blue}
\setcounter{secnumdepth}{0}
\setlength{\parindent}{0pt}
\setlength{\parskip}{5pt}
\setlength{\emergencystretch}{3em}
\newcommand{\sref}[2]{\hyperlink{source:#1:#2}{[S#2]}}
\newcommand{\cataloguescope}[1]{\paragraph{Recorded target.}#1}
\newcommand{\R}{\mathbb R}
\newcommand{\C}{\mathbb C}
\newcommand{\N}{\mathbb N}
\newcommand{\Z}{\mathbb Z}
\newcommand{\Q}{\mathbb Q}
\newcommand{\F}{\mathbb F}
\newcommand{\eps}{\varepsilon}
\begin{document}
\section{Local Well-Posedness of Lovelock Equations}\label{statement}

\subsection{Definitions and mathematical statement}

\paragraph{Shared notation.}$\mathbb N=\{1,2,\ldots\}$; the standard number fields have their usual meanings. All additional conventions are stated locally.


\paragraph{Definitions and conventions.}
Lovelock gravity uses Euler--Lagrange equations from a sum of dimensionally extended curvature densities; they remain second order in the metric despite involving higher powers of curvature. A local Cauchy theory consists of existence, uniqueness up to the relevant diffeomorphism gauge and continuous dependence for constraint-compatible initial data in specified function spaces. Hyperbolicity and the admissible branch of the equations are essential to such a theory.

\paragraph{Target retained from the canonical source catalogue.}
\paragraph{Catalogue wording retained for traceability.}
Can one establish local existence and uniqueness of solutions, together with continuous dependence on initial data, for the Lovelock-theory equations?

\paragraph{Source scope and qualification.}
The spacetime dimension, coupling constants, branch and data regime are not fixed by the extracted sentence. They must be specified in a theorem; local well-posedness on one branch cannot be presented as a result for all Lovelock data.
\subsection{Short English statement}
Establish a locally well-posed initial-value theory for Lovelock gravitational equations.
\subsection{Sources}
\begin{itemize}
\item[\hypertarget{source:statement:1}{S1}] ULAM.AI and the UnsolvedMath Contributors, UnsolvedMath, catalogue entry. Catalogue curation under CC BY 4.0; definitions and English exposition added locally.
\url{https://huggingface.co/datasets/ulamai/UnsolvedMath}.
\item[\hypertarget{source:statement:2}{S2}] Mathematical Aspects of General Relativity (2015), Oberwolfach Reports, DOI 10.4171/owr/2015/33, pp. 1867--1935. Primary report accessed 2026-10-11; relevant PDF pages 38.
\url{https://doi.org/10.4171/owr/2015/33}.
\item[\hypertarget{source:statement:3}{S3}] Papallo--Reall, local well-posedness of Lovelock theories (catalogue-listed context; not independently consulted in this preparation).
\url{https://arxiv.org/abs/1705.04370}.
\item[\hypertarget{source:statement:4}{S4}] Kovacs--Reall, well-posed Lovelock formulation (catalogue-listed context; not independently consulted in this preparation).
\url{https://arxiv.org/abs/2003.08398}.
\end{itemize}
\label{end:statement}
\end{document}
